\subsection{Two-sided density moments and antiferromagnetic mean
field}\label{proof:positive_structure:two-sided-density-moments-and-antiferromagnetic-mean-field}

The earlier density-floor criterion can be strengthened: pointwise
density comparison is unnecessary if both directions of comparison have
suitable moments.

\textbf{Theorem P5 (two-sided moment transfer).} Let \(\mu _n\) have
full support, with \(r_n=\mathrm d\mu_n/\mathrm dU_n\). Suppose that, uniformly in
dimension,

\[
\|r_n\|_{L^p(U_n)}\le M,
\qquad \|r_n^{-1}\|_{L^s(\mu_n)}\le N,
\qquad p,s\in(1,\infty].
\]

Write \(q=p/(p-1)\) and \(t=s/(s-1)\), using conjugate exponent one when
\(p\) or \(s\) is infinite. Then

\[
\boxed{\max_i I_i^{\mu_n}(f)
\ge c_0^tN^{-t}\big[M^{-q}4^{1-q}\big]^{5t/2}
\frac{v^{5qt/2}}{d^{6t}}.}
\tag{P5}
\]

The support assumption is necessary for this argument: finiteness of an
inverse-density moment measured under \(\mu\) alone would not detect
missing uniform mass at points where \(r=0\).

\textbf{Proof.} Let \(a_0=\mathbb{E}_Uf\), \(V=\operatorname{Var}_Uf\),
and \(g=(f-a_0)^2\le 4\). Hölder and the bound on \(g\) give

\[
v\le \mathbb{E}_U(rg)\le M(\mathbb{E}_Ug^q)^{1/q}
\le M4^{1-1/q}V^{1/q},
\]

and hence \(V\ge M^{-q}4^{1-q}v^q\). For \(p=\infty\), this is simply
\(v\le MV\).

For each coordinate put \(h=(D_i f)^2\), which lies in \([0,1]\). Full
support and Hölder's inequality after reversing the change of measure
give

\[
I_i^U(f)=\mathbb{E}_\mu(r^{-1}h)
\le N(\mathbb{E}_\mu h^t)^{1/t}
\le N I_i^\mu(f)^{1/t}.
\]

When \(s=\infty\), the last statement is the ordinary density-floor
comparison. Therefore \(I_i^\mu  \ge  (I_i^U/N)^t\). Choose the
coordinate supplied by UAA and substitute the variance bound.
\qed

\subsubsection{Antiferromagnetic Curie--Weiss at every fixed
strength}\label{proof:positive_structure:antiferromagnetic-curieweiss-at-every-fixed-strength}

Let \(b>0\) be fixed and define

\[
\mu_{n,-b}(x)=\frac{1}{2^nZ_{n,-b}}
\exp\left[-\frac b{2n}\left(\sum_i x_i\right)^2\right],
\qquad
Z_{n,-b}=\mathbb{E}_U e^{-bS^2/(2n)}.
\]

Since \(\mathbb{E}_US^2=n\), Jensen gives \(e^{-b/2}\le Z_{n,-b}\le 1\).
Thus its density satisfies \(\lVert r\rVert_\infty \le e^{b/2}\). For
every \(s>1\) with \((s-1)b<1\),

\[
\mathbb{E}_\mu r^{-s}
=Z_{n,-b}^{s-1}\mathbb{E}_Ue^{(s-1)bS^2/(2n)}
\le[1-(s-1)b]^{-1/2}.
\tag{5.1}
\]

For completeness, the exponential-moment inequality used here follows by
introducing a standard Gaussian \(G\):

\[
\mathbb{E}_Ue^{aS^2/(2n)}
=\mathbb{E}_G\cosh(\sqrt{a/n}\,G)^n
\le \mathbb{E}_Ge^{aG^2/2}=(1-a)^{-1/2}
\quad(0\le a<1).
\]

Choose \(s=1+1/(2b)\). Then \(t=1+2b\), \(N\le 2^{1/(2s)}\), and
\(p=\infty ,q=1\). P5 gives the explicit dimension-independent bound

\[
\boxed{\max_i I_i^{\mu_{n,-b}}(f)
\ge c_0^{1+2b}\,2^{-b}
e^{-5b(1+2b)/4}
\frac{v^{(5/2)(1+2b)}}{d^{6(1+2b)}}.}
\tag{5.2}
\]

In particular, \textbf{every fixed antiferromagnetic mean-field strength
satisfies full polynomial AA}. Constants and exponents are allowed to
depend on that fixed strength. This does not assert a uniform theorem
for arbitrary growing \(b=b_n\).

Combining (5.2) with the ferromagnetic and critical results in the
original investigation yields an exact signed-coupling classification
for

\[
\mu_{n,\lambda}(x)\propto e^{\lambda S(x)^2/(2n)},
\qquad\lambda\in\mathbb R\text{ fixed}:
\]

\[
\boxed{\text{influence AA holds if and only if }\lambda\le1.}
\]

The positive endpoint \(\lambda =1\) is included; the obstruction for
\(\lambda >1\) is the constant-variance magnetization witness. Heat-bath
influence has the same fixed-parameter classification, because its
conditional variance factor is between
\(\operatorname{sech}^2(|\lambda |)\) and one.

\subsubsection{Uniform-preserving channels preserve the moment
criterion}\label{proof:positive_structure:uniform-preserving-channels-preserve-the-moment-criterion}

Let a Markov channel preserve uniform measure. Its action on densities
can be written as a positive averaging operator \(K\) with \(K1=1\) and
\(\mathbb{E}_U Kh=\mathbb{E}_Uh\). If the input density is \(r\), the
output density is \(r_{\mathrm{out}}=Kr\). Jensen gives

\[
\|Kr\|_{L^p(U)}\le\|r\|_{L^p(U)},
\qquad \mathbb{E}_U(Kr)^{1-s}\le \mathbb{E}_Ur^{1-s}
\quad(p>1,\ s>1),
\]

because both \(z^p\) and \(z^{1-s}\) are convex. The second inequality
is precisely

\(\mathbb{E}_{\mu _{\mathrm{out}}} r_{\mathrm{out}}^{-s} \le  \mathbb{E}_\mu  r^{-s}\).

The infinite-exponent cases follow by preserving the maximum and minimum
bounds. Positive input density gives positive output density. Thus every
such channel preserves the hypotheses and constants of P5. This
conclusion concerns this sufficient moment criterion; it does not assert
that arbitrary channels preserve AA without these hypotheses.

In particular, the product spin channel \(\Pr (X_i=x\mid S_i=s)=(1+cxs)/2\),
for any \(c \in [-1,1]\), preserves uniform measure. Therefore the
antiferromagnetic mean-field bound (5.2) holds after this channel, with
constants independent of \(c\).

Applied to the measured thermal law of the commuting quantum interaction
\(+(u\cdot \sigma )\otimes (u\cdot \sigma )\) on normalized complete
graphs, the underlying Ising law is antiferromagnetic and the
measurement channel has \(c=u_z\). Consequently \textbf{all fixed
inverse temperatures and all orientations satisfy influence AA in this
antiferromagnetic subfamily}, uniformly over orientation at each fixed
inverse temperature. The bound for heat-bath influence follows as well:
conditioning through the underlying spins gives
\(|\mathbb{E}(X_i\mid X_{-i})|\le |c|\tanh  b\), so its variance factor is at
least \(\operatorname{sech}^2 b\).

\subsubsection{Positive semidefinite quadratic penalties with bounded
trace}\label{proof:positive_structure:positive-semidefinite-quadratic-penalties-with-bounded-trace}

There is a useful matrix-valued version. Let \(A_n\) be positive
semidefinite and suppose \(\operatorname{tr}(A_n)\le B\), for fixed
\(B>0\), uniformly in dimension. Absorb inverse temperature into
\(A_n\), and consider

\[
\mu_{A_n}(x)\propto \exp(-x^TA_nx/2).
\]

Its normalizer relative to uniform is at least \(e^{-B/2}\) by Jensen,
and its density is at most \(e^{B/2}\). For
\(0<\theta <1/\lVert A_n\rVert_{\mathrm{op}}\), a centered Gaussian
\(G\) of covariance \(A_n\) gives

\[
\begin{aligned}
\mathbb{E}_Ue^{\theta X^TA_nX/2}
&=\mathbb{E}_G\prod_i\cosh(\sqrt\theta\,G_i)\\
&\le \mathbb{E}_Ge^{\theta\|G\|_2^2/2}
=\det(I-\theta A_n)^{-1/2}\\
&\le\exp\left(\frac{\theta\operatorname{tr}A_n}
{2(1-\theta\|A_n\|_{\rm op})}\right).
\end{aligned}
\]

The last inequality follows by applying \(-\log (1-z)\le z/(1-z_{\max} )\)
to the nonnegative eigenvalues scaled by \(\theta\). Since
\(\lVert A_n\rVert_{\mathrm{op}}\le B\), choosing \(\theta =1/(2B)\)
bounds this moment by \(e^{1/2}\). With \(s=1+\theta\) and \(t=1+2B\),
the inverse-density calculation of Section~\ref{proof:positive_structure:antiferromagnetic-curieweiss-at-every-fixed-strength} therefore gives
\(N\le e^{1/(2s)}\). P5 implies

\[
\boxed{\max_i I_i^{\mu_{A_n}}(f)
\ge c_0^{1+2B}e^{-B-5B(1+2B)/4}
\frac{v^{(5/2)(1+2B)}}{d^{6(1+2B)}}.}
\tag{5.3}
\]

Thus arbitrary fixed-rank mean-field quadratic penalties of bounded
total spectral weight satisfy AA; rank itself need not be bounded if the
trace is. The zero matrix gives uniform measure directly. A
uniform-preserving channel preserves this conclusion by Section~\ref{proof:positive_structure:uniform-preserving-channels-preserve-the-moment-criterion}.

Equivalently, the classical Hamiltonian must admit such a positive
semidefinite representation with \textbf{uniformly bounded trace}, up to
a constant on the cube. Diagonal changes are permitted in finding a
representation because they change only that constant. Positive
semidefiniteness alone is insufficient: adding a large diagonal matrix
can make the trace grow with dimension, violating the hypothesis.

This condition distinguishes the isospectral pair in
Appendix~\ref{proof:local_interactions}. For one ferromagnetic edge of unit
strength, the two-by-two penalty matrix
\(\bigl(\begin{smallmatrix}1&-1\\-1&1\end{smallmatrix}\bigr)\) is
positive semidefinite of trace two and produces \(-x_1x_2\) up to a
constant. In contrast, for the delocalized normalized ferromagnetic
matrix with \(J 1=1\), a penalty representation has off-diagonal entries
\(-J_{ij}\). If its diagonal is \(d_i\), positivity on the all-ones vector
forces

\[
0\le\mathbf1^TA\mathbf1=\sum_i d_i-\sum_{i,j}J_{ij}
=\operatorname{tr}A-n.
\]

Every such positive semidefinite completion has trace at least \(n\). It
cannot meet the bounded-trace hypothesis, despite the bounded rank and
norm of the original interaction matrix.
